% Part: methods
% Chapter: induction
% Section: introduction

\documentclass[../../../include/open-logic-section]{subfiles}

\begin{document}

\olfileid{mth}{ind}{int}

\olsection{جان پہچان}

استقرا ثبوت دا اک ضروری طریقہ اے جیہڑا وکھ وکھ صورتاں وچ منطق، نظری کمپیوٹر سائنس، تے ریاضی دے لگ بھگ سارے شعبیاں وچ ورتیا جاندا اے۔ منطق دے کئی نتیجے ثابت کرن لئی ایہدی لوڑ اے۔

استقرا نوں اکثر استخراج دے مقابلے وچ رکھیا جاندا اے تے خاص صورتاں توں عام دعوے تک پہنچن والا استنباط دسیا جاندا اے۔ مثلاً جے اسیں کئی ہرے زمرد ویکھیے تے کوئی ایہو جیہی شے نہ ویکھیے جیہنوں زمرد آکھیے پر اوہ ہری نہ ہووے، تاں اسیں نتیجہ کڈھ سکدے آں کہ سارے زمرد ہرے نیں۔ ایہ استقرائی استنباط اے، کیوں جے ایہ کئی خاص صورتاں (ایہ زمرد ہرا اے، اوہ زمرد ہرا اے، وغیرہ) توں اک عام دعوے (سارے زمرد ہرے نیں) تک جاندا اے۔ \emph{ریاضیاتی} استقرا وی اوہ استنباط اے جیہدا نتیجہ عام دعویٰ اے، پر ایہ ایس «سادے استقرا» توں بڑی وکھری قسم دا اے۔

موٹی گل ایہ اے کہ ریاضی وچ استقرائی ثبوت ایہ نتیجہ کڈھدا اے کہ اک خاص قسم دیاں ساریاں ریاضیاتی شےواں وچ کوئی خاصیت اے۔ سبھ توں سوکھی صورت وچ استقرائی ثبوت جیہناں شےواں بارے اے، اوہ قدرتی عدد نیں۔ ایس صورت وچ استقرائی ثبوت ایہ ثابت کرن لئی ورتیا جاندا اے کہ سارے قدرتی عدداں وچ کوئی خاصیت اے۔ ایہ کم ایویں ہُندا اے:
\begin{enumerate}
    \item $0$ وچ ایہ خاصیت اے، تے
    \item جدوں وی کسے عدد~$k$ وچ ایہ خاصیت اے، تاں~$k+1$ وچ وی اے۔
\end{enumerate}
قدرتی عدداں اُتے استقرا نوں اوہناں ریاضیاتی شےواں بارے عام دعوے ثابت کرن لئی وی اکثر ورت سکدے آں جیہناں نوں عدد دتے جا سکدے نیں۔ مثلاً ہر محدود سیٹ دے رکناں دی محدود گنتی~$n$ اے۔ جے اسیں استقرا نال وکھا سکیے کہ ہر عدد~$n$ وچ ایہ خاصیت اے کہ «وڈائی~$n$ والے سارے محدود سیٹ \dots»، تاں اسیں سارے محدود سیٹاں بارے کوئی گل وکھا دتی ہووے گی۔

استقرا نوں اوہناں ریاضیاتی شےواں تک وی ودھایا جا سکدا اے جیہڑیاں \emph{استقرائی طریقے نال تعریف} کیتیاں گئیاں نیں۔ مثلاً رسمی زبان دیاں عبارتاں، جیویں پہلے درجے دی منطق دیاں عبارتاں، استقرائی طریقے نال تعریف ہُندیاں نیں۔ \emph{ساختی استقرا} ایہو جیہیاں ساریاں عبارتاں بارے نتیجے ثابت کرن دا اک طریقہ اے۔ خاص کر کے ساختی استقرا منطق وچ بڑا کم آؤندا اے تے بہت ورتیا جاندا اے۔

\end{document}
